
% ---- fig033_retta11.pgf ----
% (c) 2014 Daniele Zambelli - daniele.zambelli@gmail.com

%%%
% Retta decrescente zero in 3
%%%%
 
\begin{tikzpicture}[x=1.5mm, y=1.5mm, smooth]

\input{\folder lbr/lib_rettadec.pgf}
\node [above] {$3$};

\end{tikzpicture}
% -----------------

% ---- fig034_retta12.pgf ----
% (c) 2014 Daniele Zambelli - daniele.zambelli@gmail.com

%%%
% Retta crescente zero in -2
%%%%
 
\begin{tikzpicture}[x=1.5mm, y=1.5mm, smooth]

\input{\folder lbr/lib_rettacre.pgf}
\node [above] {$-2$};

\end{tikzpicture}
% -----------------

% ---- fig035_grafosegni03.pgf ----
% (c) 2014 Daniele Zambelli - daniele.zambelli@gmail.com

%%%
% Studio dei segni di una frazione
%%%%
 
\begin{tikzpicture}[x=2.5mm, y=5mm, smooth]

\coordinate (a_top) at (-3.3, 1);
\coordinate (a_bottom) at (-3.3, -2);
\coordinate (b_top) at (3.3, 1);
\coordinate (b_bottom) at (3.3, -2);

\input{\folder lbr/lib_grafo3assi.pgf}
\node [above] at (-3.3, 1) {$-2$};
\node [above] at (3.3, 1) {$+3$};

\node [above left] at (-10, 0) {$-3 x + 4$};
\node [above] at (-6.5, 0) {$+$};
\node [above] at (0, 0) {$+$};
\draw (3.3, .5) circle (3pt);
\node [above] at (6.5, 0) {$-$};

\node [above left] at (-10, -1) {$x + 2$};
\node [above] at (-6.5, -1) {$-$};
\draw (-3.3 -.4, -.5 -.2) -- (-3.3 +.4, -.5 +.2) 
      (-3.3 -.4, -.5 +.2) -- (-3.3 +.4, -.5 -.2);
\node [above] at (0, -1) {$+$};
\node [above] at (6.5, -1) {$+$};

\node [above left] at (-10, -2.15) {$f(x)$};
\node [above] at (-6.5, -2) {$-$};
\draw (-3.3 -.4, -1.5 -.2) -- (-3.3 +.4, -1.5 +.2) 
      (-3.3 -.4, -1.5 +.2) -- (-3.3 +.4, -1.5 -.2);
\node [above] at (0, -2) {$+$};
\draw (3.3, -1.5) circle (3pt);
\node [above] at (6.5, -2) {$-$};

\end{tikzpicture}
% -----------------

% ---- fig036_diseq02.pgf ----
% (c) 2014 Daniele Zambelli - daniele.zambelli@gmail.com

%%%
% Valori esterni all'intervallo -2; 3
%%%%
 
\begin{tikzpicture}[x=1.5mm, y=1.5mm, smooth]

% \clip (-7.5, -5.5) rectangle (10.9, 10.9);

\coordinate (m_i) at (-10, 0);
\coordinate (a) at (-3, 0);
\coordinate (b) at (3, 0);
\coordinate (p_i) at (10, 0);

\input{\magdir assiepiani/asse10x.pgf}

\begin{scope}[blue,thick]
\draw [-,decorate,decoration=snake] (m_i) -- (a);
\draw[fill=white] (a) circle (2pt) node [above] {$-2$};
\draw [-,decorate,decoration=snake] (b) -- (p_i);
\draw[fill] (b) circle (2pt) node [above] {$3$};
\end{scope}

\end{tikzpicture}
% -----------------

% ---- fig037_grafosegni031.pgf ----
% (c) 2014 Daniele Zambelli - daniele.zambelli@gmail.com

%%%
% Segni e soluzione di una disequazione fratta
%%%%
 
\begin{tikzpicture}[x=2.5mm, y=5mm, smooth]

\coordinate (m_i) at (-10, -2);
\coordinate (a_top) at (-3.3, 1);
\coordinate (a_bottom) at (-3.3, -2);
\coordinate (b_top) at (3.3, 1);
\coordinate (b_bottom) at (3.3, -2);
\coordinate (p_i) at (10, -2);

\input{\folder lbr/lib_grafo3assi.pgf}
\node [above] at (-3.3, 1) {$-2$};
\node [above] at (3.3, 1) {$+3$};

\node [above left] at (-10, 0) {$-3 x + 4$};
\node [above] at (-6.5, 0) {$+$};
\node [above] at (0, 0) {$+$};
\draw (3.3, .5) circle (3pt);
\node [above] at (6.5, 0) {$-$};

\node [above left] at (-10, -1) {$x + 2$};
\node [above] at (-6.5, -1) {$-$};
\draw (-3.3 -.4, -.5 -.2) -- (-3.3 +.4, -.5 +.2) 
      (-3.3 -.4, -.5 +.2) -- (-3.3 +.4, -.5 -.2);
\node [above] at (0, -1) {$+$};
\node [above] at (6.5, -1) {$+$};

\node [above left] at (-10, -2.15) {$f(x)$};
\node [above] at (-6.5, -2) {$-$};
\draw (-3.3 -.4, -1.5 -.2) -- (-3.3 +.4, -1.5 +.2) 
      (-3.3 -.4, -1.5 +.2) -- (-3.3 +.4, -1.5 -.2);
\node [above] at (0, -2) {$+$};
\draw (3.3, -1.5) circle (3pt);
\node [above] at (6.5, -2) {$-$};

\begin{scope}[blue,thick]
\draw [-,decorate,decoration=snake] (m_i) -- (a_bottom);
\draw[fill=white] (a_bottom) circle (2pt);
\draw [-,decorate,decoration=snake] (b_bottom) -- (p_i);
\draw[fill] (b_bottom) circle (2pt);
\end{scope}

\end{tikzpicture}
% -----------------

% ---- fig038_cerchietto.pgf ----
% (c) 2014 Daniele Zambelli - daniele.zambelli@gmail.com

%%%
% Cerchietto
%%%%
 
\begin{tikzpicture}[smooth]

\draw (0, 0) circle (3pt);

\end{tikzpicture}
% -----------------

% ---- fig039_crocetta.pgf ----
% (c) 2014 Daniele Zambelli - daniele.zambelli@gmail.com

%%%
% Studio dei segni di una frazione
%%%%
 
\begin{tikzpicture}[smooth]

\draw ( -.1,  -.1) -- (+.1, +.1) 
      ( -.1,  +.1) -- (+.1, -.1);

\end{tikzpicture}
% -----------------

% ---- fig040_escluso.pgf ----
% (c) 2014 Daniele Zambelli - daniele.zambelli@gmail.com

%%%
% Studio dei segni di una frazione
%%%%
 
\begin{tikzpicture}[smooth]

\draw[blue, thick, fill=white] (0, 0) circle (2pt);

\end{tikzpicture}
% -----------------

% ---- fig041_incluso.pgf ----
% (c) 2014 Daniele Zambelli - daniele.zambelli@gmail.com

%%%
% Studio dei segni di una frazione
%%%%
 
\begin{tikzpicture}[smooth]

\draw[fill, blue, thick] (0, 0) circle (2pt);

\end{tikzpicture}
% -----------------

% ---- fig053_diseq06.pgf ----
% (c) 2014 Daniele Zambelli - daniele.zambelli@gmail.com

%%%
% Intervallo ]oo; 2/9[ 
%%%%
 
\begin{tikzpicture}[x=1.5mm, y=1.5mm, smooth]

% \clip (-7.5, -5.5) rectangle (10.9, 10.9);

\coordinate (a) at (-10, 0);
\coordinate (b) at (+10, 0);

\input{\magdir assiepiani/asse10x.pgf}

\end{tikzpicture}
% -----------------

% ---- fig106_ret.pgf ----
% (c) 2012 Dimitrios Vrettos - d.vrettos@gmail.com
\begin{tikzpicture}[font=\small,x=10mm, y=5mm]

\draw[->] (0,0) -- (8,0) node [below right] () {$r$};
\node[above]  at (4,0) {0};
\begin{scope}[blue,thick,->]
\draw (4,0) -- (8,0);
\draw[fill=white] (4,0)circle (1.5pt);
\end{scope}

\end{tikzpicture}% -----------------

% ---- fig108_eser.pgf ----
% (c) 2012 Dimitrios Vrettos - d.vrettos@gmail.com
\begin{tikzpicture}[font=\small,x=10mm, y=5mm]

\draw[->] (0,0) -- (8,0) node [below right] () {$r$};
\node[above]  at (4,0) {$-3$};
\begin{scope}[blue,thick]
\draw (0,0) -- (4,0);
\draw[fill=white] (4,0)circle (1.5pt);
\end{scope}

\end{tikzpicture}% -----------------

% ---- fig109_eser.pgf ----
% (c) 2012 Dimitrios Vrettos - d.vrettos@gmail.com
\begin{tikzpicture}[font=\small,x=10mm, y=5mm]

\draw[->] (0,0) -- (8,0) node [below right] () {$r$};
\node[above]  at (4,0) {$2$};
\begin{scope}[blue,thick,->]
\draw (4,0) -- (8,0);
\draw[fill=blue] (4,0)circle (1.5pt);
\end{scope}

\end{tikzpicture}% -----------------

% ---- fig110_eser.pgf ----
% (c) 2012 Dimitrios Vrettos - d.vrettos@gmail.com
\begin{tikzpicture}[font=\small,x=10mm, y=5mm]

\draw[->] (0,0) -- (8,0) node [below right] () {$r$};
\node[above]  at (2,0) {$-2$};
\node[above]  at (6,0) {2};
\begin{scope}[blue,thick]
\draw (2,0) -- (6,0);
\foreach \x in {2,6}
\draw[fill=white] (\x,0)circle (1.5pt);
\end{scope}

\end{tikzpicture}% -----------------

% ---- fig111_eser.pgf ----
% (c) 2012 Dimitrios Vrettos - d.vrettos@gmail.com
\begin{tikzpicture}[font=\small,x=10mm, y=5mm]

\draw[->] (0,0) -- (8,0) node [below right] () {$r$};
\node[above]  at (2,0) {$3$};
\node[above]  at (6,0) {5};
\begin{scope}[blue,thick]
\draw (2,0) -- (6,0);
\draw[fill=white] (2,0)circle (1.5pt);
\draw[fill=blue] (6,0)circle (1.5pt);
\end{scope}

\end{tikzpicture}% -----------------

% ---- fig112_eser.pgf ----
% (c) 2012 Dimitrios Vrettos - d.vrettos@gmail.com
\begin{tikzpicture}[font=\small,x=10mm, y=5mm]

\draw[->] (0,0) -- (8,0) node [below right] () {$r$};
\node[above]  at (2,0) {$-1$};
\node[above]  at (6,0) {0};
\begin{scope}[blue,thick]
\draw (2,0) -- (6,0);
\draw[fill=blue] (2,0)circle (1.5pt);
\draw[fill=blue] (6,0)circle (1.5pt);
\end{scope}

\end{tikzpicture}% -----------------

% ---- fig113_eser.pgf ----
% (c) 2012 Dimitrios Vrettos - d.vrettos@gmail.com
\begin{tikzpicture}[font=\small,x=10mm, y=5mm]

\draw[->] (0,0) -- (8,0) node [below right] () {$r$};
\node[above]  at (2,0) {$1$};
\node[above]  at (6,0) {2};
\begin{scope}[blue,thick]
\draw (2,0) -- (6,0);
\draw[fill=blue] (2,0)circle (1.5pt);
\draw[fill=white] (6,0)circle (1.5pt);
\end{scope}

\end{tikzpicture}% -----------------

% ---- lib_grafo3assi.pgf ----
% (c) 2014 Daniele Zambelli - daniele.zambelli@gmail.com

%%%
% Grafo per il calcolo del segno con tre assi
%%%%
 
\input{\magdir assiepiani/asse10x.pgf}
\begin{scope}[yshift= -.5cm]
  \input{\magdir assiepiani/asse10x.pgf}
  \begin{scope}[yshift= -.5cm]
    \input{\magdir assiepiani/asse10x.pgf}
  \end{scope}
\end{scope}

\draw [-] [] (a_top) -- (a_bottom);
\draw [-] [] (b_top) -- (b_bottom);
% -----------------

% ---- lib_grafo4assi.pgf ----
% (c) 2014 Daniele Zambelli - daniele.zambelli@gmail.com

%%%
% Grafo per il calcolo del segno con tre assi
%%%%
 
\input{\magdir assiepiani/asse10x.pgf}
\begin{scope}[yshift= -.5cm]
  \input{\magdir assiepiani/asse10x.pgf}
  \begin{scope}[yshift= -.5cm]
    \input{\magdir assiepiani/asse10x.pgf}
    \begin{scope}[yshift= -.5cm]
      \input{\magdir assiepiani/asse10x.pgf}
    \end{scope}
  \end{scope}
\end{scope}

\draw [-] [] (a_top) -- (a_bottom);
\draw [-] [] (b_top) -- (b_bottom);
\draw [-] [] (c_top) -- (c_bottom);
% -----------------

% ---- lib_grafo6assi.pgf ----
% (c) 2014 Daniele Zambelli - daniele.zambelli@gmail.com

%%%
% Grafo per il calcolo del segno con sei assi
%%%%
 
\input{\magdir assiepiani/asse10x.pgf}
\begin{scope}[yshift= -.5cm]
  \input{\magdir assiepiani/asse10x.pgf}
  \begin{scope}[yshift= -.5cm]
    \input{\magdir assiepiani/asse10x.pgf}
    \begin{scope}[yshift= -.5cm]
      \input{\magdir assiepiani/asse10x.pgf}
      \begin{scope}[yshift= -.5cm]
        \input{\magdir assiepiani/asse10x.pgf}
        \begin{scope}[yshift= -.5cm]
          \input{\magdir assiepiani/asse10x.pgf}
        \end{scope}
      \end{scope}
    \end{scope}
  \end{scope}
\end{scope}

\draw [-] [] (a_top) -- (a_bottom);
\draw [-] [] (b_top) -- (b_bottom);
\draw [-] [] (c_top) -- (c_bottom);
\draw [-] [] (d_top) -- (d_bottom);
\draw [-] [] (e_top) -- (e_bottom);
% -----------------

% ---- lib_rettacre.pgf ----
% (c) 2014 Daniele Zambelli - daniele.zambelli@gmail.com

%%%
% Retta crescente con segni
%%%%
 
\coordinate (inizio) at (-10, -4);
\coordinate (zero) at (0, 0);
\coordinate (fine) at (10, 4);

\input{\magdir assiepiani/asse10x.pgf}

\draw [-] [ultra thick, red!50!black] (inizio) -- (zero);
\draw [-] [ultra thick, blue!50!black] (zero) -- (fine);

\node [xshift=-25, yshift=-3, above] at (zero) {$-$};
\draw[blue, thick, fill=white] (zero) circle (2pt);
\node [xshift=25, yshift=-3, above] at (zero) {$+$};
% -----------------

% ---- lib_rettacrek.pgf ----
% (c) 2014 Daniele Zambelli - daniele.zambelli@gmail.com

%%%
% Retta k crescente con segni
%%%%
 
\coordinate (inizio) at (-10, -4);
\coordinate (zero) at (0, 0);
\coordinate (fine) at (10, 4);

\input{\magdir assiepiani/asse10k.pgf}

\draw [-] [ultra thick, red!50!black] (inizio) -- (zero);
\draw [-] [ultra thick, blue!50!black] (zero) -- (fine);

\node [xshift=-25, yshift=-3, above] at (zero) {$-$};
\draw[blue, thick, fill=white] (zero) circle (2pt);
\node [xshift=25, yshift=-3, above] at (zero) {$+$};
% -----------------

% ---- lib_rettadec.pgf ----
% (c) 2014 Daniele Zambelli - daniele.zambelli@gmail.com

%%%
% Retta decrescente con segni
%%%%
 
\coordinate (inizio) at (-10, 4);
\coordinate (zero) at (0, 0);
\coordinate (fine) at (10, -4);

\input{\magdir assiepiani/asse10x.pgf}

\draw [-] [ultra thick, blue!50!black] (inizio) -- (zero);
\draw [-] [ultra thick, red!50!black] (zero) -- (fine);

\node [xshift=-25, yshift=-3, above] at (zero) {$+$};
\draw[blue, thick, fill=white] (zero) circle (2pt);
\node [xshift=25, yshift=-3, above] at (zero) {$-$};
% -----------------

% ---- lib_rettaneg.pgf ----
% (c) 2014 Daniele Zambelli - daniele.zambelli@gmail.com

%%%
% Retta crescente
%%%%
 
\coordinate (inizio) at (-10, -3);
\coordinate (zero) at (0, 0);
\coordinate (fine) at (10, -3);

\input{\magdir assiepiani/asse10x.pgf}

\draw [-] [ultra thick, red!50!black] (inizio) -- (fine);

\node [xshift=-25, yshift=-3, above] at (zero) {$-$};
\node [xshift=25, yshift=-3, above] at (zero) {$-$};
% -----------------

% ---- lib_rettapos.pgf ----
% (c) 2014 Daniele Zambelli - daniele.zambelli@gmail.com

%%%
% Segno di una costante positiva.
%%%%
 
\coordinate (inizio) at (-10, 3);
\coordinate (zero) at (0, 0);
\coordinate (fine) at (10, 3);

\input{\magdir assiepiani/asse10x.pgf}

\draw [-] [ultra thick, blue!50!black] (inizio) -- (fine);

\node [xshift=-25, yshift=-3, above] at (zero) {$+$};
\node [xshift=25, yshift=-3, above] at (zero) {$+$};
% -----------------
